\section{Maximum Depth (Height) of a Binary Tree}
\label{sec:trees-max-depth}

\problemstatement{Given the root of a binary tree, return its maximum
depth: the number of nodes along the longest path from the root down to a
leaf.}

\begin{patternbox}
\emph{``Height,'' ``depth,'' ``longest path to a leaf''} is the archetypal
\textbf{bottom-up} recursion, and the first concrete instance of this
chapter's core rhythm. A node's height depends on its children's heights,
so you must have both child answers \emph{before} you can answer for the
node -- that ordering \emph{is} postorder.
\end{patternbox}

\subsection*{Understanding the problem carefully}

The height of a tree rooted at a node is $1$ (for the node itself) plus
the height of its taller subtree. An empty tree has height $0$. This is
the purest example of ``ask both children, combine'': the combine step is
simply $1+\max(\text{left height}, \text{right height})$.

\begin{ideabox}[Height = 1 + max of the Two Child Heights]
Recurse to the left and right children, obtain their heights, and return
one more than the larger. The base case (null returns $0$) is what
terminates the recursion and seeds the count. No global state is needed:
the answer flows straight up the return values.
\end{ideabox}

\subsection*{Algorithm}

\begin{enumerate}
  \item \textsc{Height}(\textit{node}):
  \begin{enumerate}
    \item If \textit{node} is null, return $0$.
    \item $L = $ \textsc{Height}(\textit{node}.left); $R = $
          \textsc{Height}(\textit{node}.right).
    \item Return $1 + \max(L, R)$.
  \end{enumerate}
\end{enumerate}

\subsection*{Dry run}

\dryrun{Tree: root $1$; left $2$ (leaf); right $3$ whose right child is
$4$ (leaf).}

\begin{itemize}
  \item Height($2$)$=1$ (both children null $\to 1+\max(0,0)$).
  \item Height($4$)$=1$. Height($3$)$=1+\max(0,1)=2$.
  \item Height($1$)$=1+\max(1,2)=3$.
\end{itemize}

\subsection*{C++ implementation}

\begin{lstlisting}
#include <bits/stdc++.h>
using namespace std;

struct TreeNode {
    int val;
    TreeNode* left;
    TreeNode* right;
    TreeNode(int v) : val(v), left(nullptr), right(nullptr) {}
};

int maxDepth(TreeNode* node) {
    if (node == nullptr) return 0;
    int L = maxDepth(node->left);
    int R = maxDepth(node->right);
    return 1 + max(L, R);
}

int main() {
    TreeNode* root = new TreeNode(1);
    root->left = new TreeNode(2);
    root->right = new TreeNode(3);
    root->right->right = new TreeNode(4);

    cout << "Max depth: " << maxDepth(root) << "\n";
    return 0;
}
\end{lstlisting}

\begin{complexitybox}
\textbf{Time Complexity.} $O(n)$ -- each node contributes $O(1)$ combine
work. \textbf{Space Complexity.} $O(h)$ for the recursion stack.
\end{complexitybox}

\begin{takeawaybox}
Height is the ``hello world'' of bottom-up tree recursion, and the exact
subroutine reused inside balanced-tree, diameter, and max-path-sum
problems. Each of those is ``compute height, but also check or update
something extra along the way.''
\end{takeawaybox}
